% !TeX program = pdflatex
\documentclass[11pt,a4paper]{amsart}

\newcommand{\notelanguage}{finnish}
% Shared preamble for course notes and exercise sets.

\usepackage[T1]{fontenc}
\usepackage{lmodern}
\usepackage[english]{babel}

\usepackage[a4paper,margin=30mm]{geometry}
\usepackage{microtype}
\usepackage{graphicx}
\usepackage{enumitem}

\usepackage{amsmath}
\usepackage{amssymb}
\usepackage{amsthm}
\usepackage{mathtools}
\usepackage{aliascnt}

\usepackage[hidelinks,pdfusetitle]{hyperref}

\numberwithin{equation}{section}
\setlist{itemsep=0.25em,topsep=0.5em}

\theoremstyle{plain}
\newtheorem{theorem}{Theorem}[section]
\newaliascnt{lemma}{theorem}
\newtheorem{lemma}[lemma]{Lemma}
\aliascntresetthe{lemma}
\newaliascnt{proposition}{theorem}
\newtheorem{proposition}[proposition]{Proposition}
\aliascntresetthe{proposition}
\newaliascnt{corollary}{theorem}
\newtheorem{corollary}[corollary]{Corollary}
\aliascntresetthe{corollary}

\theoremstyle{definition}
\newaliascnt{definition}{theorem}
\newtheorem{definition}[definition]{Definition}
\aliascntresetthe{definition}
\newaliascnt{example}{theorem}
\newtheorem{example}[example]{Example}
\aliascntresetthe{example}
\newaliascnt{problem}{theorem}
\newtheorem{problem}[problem]{Problem}
\aliascntresetthe{problem}

\theoremstyle{remark}
\newaliascnt{remark}{theorem}
\newtheorem{remark}[remark]{Remark}
\aliascntresetthe{remark}

\usepackage[nameinlink,noabbrev]{cleveref}

\crefname{theorem}{theorem}{theorems}
\crefname{lemma}{lemma}{lemmas}
\crefname{proposition}{proposition}{propositions}
\crefname{corollary}{corollary}{corollaries}
\crefname{definition}{definition}{definitions}
\crefname{example}{example}{examples}
\crefname{problem}{problem}{problems}
\crefname{remark}{remark}{remarks}

\DeclareMathOperator{\im}{im}
\DeclareMathOperator{\rank}{rank}
\DeclarePairedDelimiter{\abs}{\lvert}{\rvert}
\DeclarePairedDelimiter{\norm}{\lVert}{\rVert}
\DeclarePairedDelimiter{\set}{\lbrace}{\rbrace}


\title{Johdatus logiikkaan I\\Harjoituksen 3 ratkaisut}
\author{}
\date{}

\begin{document}

\exerciseheading{Johdatus logiikkaan I}{Harjoituksen 3 ratkaisut}

\begin{exercise}
  Mitkä seuraavista lauseista ovat disjunktiivisessa
  normaalimuodossa? Perustele.
\begin{enumerate}[label=\alph*), nosep]
  \item \((p_0 \land \neg p_0) \lor (\neg p_3 \land p_1)\)
  \item \((p_0 \land p_0) \lor p_0\)
  \item \(\neg p_1 \land p_2 \land p_3\)
  \item \(\neg\neg p_0 \lor p_1 \lor p_2\)
  \item \((p_0 \land \neg p_1 \land p_2) \lor \neg(p_1 \lor p_2)\)
  \item \((p_0 \lor \neg p_1) \land (p_1 \lor p_2)\)
  \item \(\neg p_1 \rightarrow (p_0 \land \neg p_1)\)
\end{enumerate}
\end{exercise}

\begin{solution}
\leavevmode\par
\begin{enumerate}[label=\alph*), nosep]
\item \(A = (p_0 \land \neg p_0)\) ja \(B = (\neg p_3 \land p_1)\) ovat
  literaalien konjunktioita, joten \(A \lor B\) on DNF.
\item \(A = (p_0 \land p_0)\) ja \(B = p_0\) ovat literaalien konjunktioita,
  joten \(A \lor B\) on DNF.
\item \(\neg p_1 \land p_2 \land p_3\) on literaalien konjunktio, joten se
  on DNF yhden disjunktin tapauksena.
\item \(\neg \neg p_0\) ei ole literaali, joten lause ei ole DNF.
\item Disjunkti \(\neg(p_1 \lor p_2)\) ei ole literaalien konjunktio, joten
  lause ei ole DNF.
\item Lause ei ole DNF, sillä se on disjunktioiden konjunktio eikä
  literaalien konjunktioiden disjunktio.
\item Pääkonnektiivi on \(\to\), joten lause ei ole DNF.
\end{enumerate}
\end{solution}

\begin{exercise}
Mitkä edellisen tehtävän lauseista ovat konjunktiivisessa normaalimuodossa?
\end{exercise}

\begin{solution}
\leavevmode\par
\begin{enumerate}[label=\alph*), nosep]
\item Lause ei ole CNF, sillä disjunktion sisällä esiintyy konjunktioita,
joten se ei ole literaalien disjunktio.
\item Lause ei ole CNF, sillä disjunktion sisällä esiintyy konjunktio,
joten se ei ole literaalien disjunktio.
\item Lause on CNF, sillä se on yhden literaalin disjunktioiden konjunktio.
\item \(\neg \neg p_0\) ei ole literaali, joten lause ei ole CNF.
\item Lause ei ole CNF, sillä disjunktion jäsenet eivät ole literaaleja.
\item Pääkonnektiivi on \(\land\), ja konjunktit \(A = (p_0 \lor \neg p_1)\)
ja \(B = (p_1 \lor p_2)\) ovat literaalien disjunktioita, joten lause
on CNF.
\item Pääkonnektiivi on \(\to\), joten lause ei ole CNF.
\end{enumerate}
\end{solution}

\begin{exercise}
Anna esimerkki propositiolauseesta, jossa esiintyy vähintään kahta eri
propositiosymbolia ja joka on
\begin{enumerate}[label=\alph*), nosep]
\item sekä disjunktiivisessa että konjunktiivisessa normaalimuodossa,
\item disjunktiivisessa, muttei konjunktiivisessa normaalimuodossa,
\item konjunktiivisessa, muttei disjunktiivisessa normaalimuodossa,
\item ei disjunktiivisessa eikä konjunktiivisessa normaalimuodossa.
\end{enumerate}
Pyri keksimään muita esimerkkejä kuin tehtävässä 1 esiintyviä.
\end{exercise}

\begin{solution}
\leavevmode\par
\begin{enumerate}[label=\alph*), nosep]
\item \(p_0 \lor p_1\) on sekä DNF että CNF.
\item \((p_0 \land p_1) \lor (p_2 \land p_3)\) on DNF, muttei CNF.
\item \((p_0 \lor p_1) \land (p_2 \lor p_3)\) on CNF, muttei DNF.
\item \(p_0 \leftrightarrow p_1\) ei ole DNF eikä CNF, sillä pääkonnektiivi
on \(\leftrightarrow\).
\end{enumerate}
\end{solution}

\begin{exercise}
Etsi disjunktiivisessa normaalimuodossa oleva lause, joka on loogisesti
ekvivalentti lauseen
\[
(p_1 \leftrightarrow (p_0 \rightarrow \neg p_2))
\land (p_0 \lor (\neg p_2 \land \neg p_1))
\]
kanssa.
\end{exercise}

\begin{solution}
Muodostamalla lauseen totuustaulu ja soveltamalla DNF-konstruktiota
saadaan seuraava ekvivalentti DNF-muotoinen lause:
\[
A = (p_0 \land \neg p_1 \land p_2) \lor (p_0 \land p_1 \land \neg p_2).
\]
\end{solution}

\begin{exercise}
Etsi disjunktiivisessa normaalimuodossa oleva lause, joka on loogisesti
ekvivalentti lauseen
\[
(p_1 \land (p_0 \rightarrow p_2))
\leftrightarrow (p_2 \rightarrow (p_1 \rightarrow p_2))
\]
kanssa.
\end{exercise}

\begin{solution}
Muodostamalla lauseen totuustaulu ja soveltamalla DNF-konstruktiota
saadaan seuraava ekvivalentti DNF-muotoinen lause:
\[
A = (\neg p_0 \land p_1 \land \neg p_2) \lor (\neg p_0 \land p_1 \land p_2)
\lor (p_0 \land p_1 \land p_2).
\]
\end{solution}

\begin{exercise}
Etsi konjunktiivisessa normaalimuodossa oleva lause, joka on loogiesti
ekvivalentti lauseen
\[
\neg(p_0 \leftrightarrow p_1) \rightarrow (p_2 \land \neg p_1)
\]
kanssa.
\end{exercise}

\begin{solution}
Olkoon \(A = \neg(p_0 \leftrightarrow p_1) \rightarrow
(p_2 \land \neg p_1)\) alkuperäinen lause.
Muodostamalla lauseen \(A\) totuustaulu ja soveltamalla DNF-konstruktiota
totuustaulun 0-riveihin saadaan seuraava lauseen \(\neg A\) kanssa
ekvivalentti DNF-muotoinen lause:
\[
\neg A \iff (\neg p_0 \land p_1 \land \neg p_2) \lor (\neg p_0 \land
p_1 \land p_2)
\lor (p_0 \land \neg p_1 \land \neg p_2).
\]
Negatoimalla saatu lause ja soveltamalla De Morganin kaavoja sekä
kaksoisnegaation lakia saadaan seuraava CNF-muotoinen lause:
\[
\neg \neg A \iff A \iff (p_0 \lor \neg p_1 \lor p_2) \land (p_0 \lor
\neg p_1 \lor \neg p_2)
\land (\neg p_0 \lor p_1 \lor p_2).
\]
\end{solution}

\begin{exercise}
Kuinka pitkä propositiolause saadaan, kun kurssin menetelmällä etsitään
annetulle propositiolauseelle \(A\) ekvivalentti DNF-muotoinen lause? Tässä
tehtävässä kuuluu miettiä myös, miten \(A\):n ja DNF-muotoisen lauseen
pituutta on mielekästä mitata. Vastaus ei välttämättä ole yksikäsitteinen.
Keskustelkaa laskuharjoitustilaisuudesta esitettyjen vastausten hyvistä ja
huonoista puolista. Löytyykö yhteinen mielipide parhaasta vastauksesta?
\end{exercise}

\begin{solution}
Olkoon \(n\) lauseessa \(A\) esiintyvien eri propositiosymbolien lukumäärä
ja \(t\) lauseen \(A\) totuustaulun tosien rivien lukumäärä.
Oletetaan ensin, että \(t \geq 1\).
DNF-konstruktiolla saadussa lauseessa on \(t\) disjunktia
ja jokaisessa näistä on \(n\) literaalia. Täten DNF-muotoisessa lauseessa
on \(nt\) literaalien esiintymää, \(t(n - 1)\) konjunktiomerkkiä
(\(\land\)) ja \(t - 1\) disjunktiomerkkiä (\(\lor\)).
Jos lasketaan jokainen literaali yhdeksi yksiköksi eikä oteta sulkeita
huomioon, saadaan pituudeksi
\[
nt + t(n - 1) + (t - 1) = 2nt - 1.
\]
Koska \(t \leq 2^n\), saadaan yläraja \(\mathcal{O}(n2^n)\). Negaatioiden
ja sulkeiden laskeminen erikseen muuttaisi tätä tarkkaa tulosta,
mutta asymptoottinen
yläraja pysyy ennallaan.

Jos \(t=0\), lause on aina epätosi, ja sen DNF-muodoksi voidaan valita
\(p_0\land\neg p_0\). Tällöin pituus on vakio.

Olkoon \(m\) lauseen \(A\) pituus laskettuna propositiosymbolien ja
konnektiivien esiintymien määränä. Koska \(n \leq m\), saadaan asymptoottinen
yläraja \(\mathcal{O}(m2^m)\).

DNF-konstruktio voi tuottaa eksponentiaalisesti pidemmän lauseen
alkuperäiseen verrattuna. Esimerkiksi lauseen
\[
A_n=(p_1\lor\neg p_1)\land\cdots\land(p_n\lor\neg p_n)
\]
pituus on verrannollinen lukuun \(n\), mutta kaikki sen totuustaulun
\(2^n\) riviä ovat tosia. Siksi konstruktio tuottaa \(2^n\) disjunktia,
joissa jokaisessa on \(n\) literaalia. Edellä käytetyllä laskutavalla
saadun DNF-lauseen pituus on siis \(2n2^n-1\).
Tämä koskee konstruktion tuottamaa lausetta: lyhyempi ekvivalentti
DNF-lause voi olla olemassa. Tässä esimerkissä sellainen on
\(p_1\lor\neg p_1\).
\end{solution}

\begin{exercise}
\textbf{Peircen nuoli} on seuraava konnektiivi:
\[
\begin{array}{cc|c}
A & B & A \mathbin{\downarrow} B \\
\hline
1 & 1 & 0 \\
1 & 0 & 0 \\
0 & 1 & 0 \\
0 & 0 & 1
\end{array}
\]
Osoita, että \(\{\downarrow\}\) on täydellinen konnektiivijoukko.
\end{exercise}

\begin{solution}
Konnektiivijoukko \(\set{\neg, \land}\) on täydellinen. Riittää siis osoittaa,
että pystymme esittämään konnektiivit \(\neg\) ja \(\land\) Peircen
nuolen avulla.

Olkoon \(f\) mielivaltainen totuusfunktio. Koska \(\set{\neg, \land}\) on
täydellinen konnektiivijoukko, on olemassa propositiolause \(A\), jolle
\(f_A=f\) ja joka sisältää ainoastaan
propositiosymboleita \(p_0, \cdots, p_n\) ja konnektiiveja \(\neg\)
ja \(\land\).

Osoitetaan induktiolla \(A\):n rakenteen suhteen, että on olemassa
sen kanssa ekvivalentti propositiolause \(A^{'}\), joka sisältää ainoastaan
propositiosymboleita \(p_0, \cdots, p_n\) ja konnektiivia \(\downarrow\).

Jos \(A = p_i, 0 \leq i \leq n\), niin valitaan \(A^{'} = p_i\).

Jos \(A = \neg B\), induktio-oletuksen nojalla on olemassa propositiolause
\(B^{'}\), jolle \(B^{'} \iff B\) ja joka sisältää vain
propositiosymboleita \(p_0, \cdots, p_n\) ja
konnektiivia \(\downarrow\). Sitten
\[
A = \neg B \iff \neg B^{'} \iff B^{'} \downarrow B^{'}.
\]
Valitaan \(A^{'} = B^{'} \downarrow B^{'}\).

Jos \(A = B \land C\), induktio-oletuksen nojalla on olemassa propositiolauseet
\(B^{'}\) ja \(C^{'}\), joille \(B^{'} \iff B\) ja \(C^{'} \iff C\) ja
jotka sisältävät vain propositiosymboleita \(p_0, \cdots, p_n\) ja
konnektiivia \(\downarrow\). Sitten
\[
\begin{aligned}
A = B \land C &\iff B^{'} \land C^{'} \\
&\iff \neg (\neg B^{'} \lor \neg C^{'}) \\
&\iff (B^{'} \downarrow B^{'}) \downarrow (C^{'} \downarrow C^{'}).
\end{aligned}
\]
Valitaan \(A^{'} = (B^{'} \downarrow B^{'}) \downarrow (C^{'}
\downarrow C^{'})\).

Koska \(f_{A^{'}} = f_{A} = f\), \(A^{'}\) määrittelee funktion \(f\)
Peircen nuolen avulla. Siis \(\set{\downarrow}\) on täydellinen
konnektiivijoukko.
\end{solution}

\begin{exercise}
Osoita, että \(\{\rightarrow\}\) ei ole täydellinen konnektiivijoukko.
(Vihje: Eräästä viime viikon tehtävästä on apua.)
\end{exercise}

\begin{solution}
Osoitetaan, että konnektiivilla \(\to\) ei voida esittää negaatiota,
eli funktiota \(f \colon \set{0, 1} \to \set{0, 1}\), jolla
\(f(x) = 1 - x\). Olkoon \(A\) propositiolause, joka
sisältää vain propositiosymbolin \(p_0\) ja konnektiivia \(\to\).
Olkoon \(v\) totuusjakauma, jolla \(v(p_0) = 1\).
Osoitetaan induktiolla \(A\):n rakenteen suhteen, että \(v(A)=1\).

Jos \(A\) on propositiosymboli, niin \(A=p_0\) ja siten \(v(A) = 1\).

Jos \(A = (B \to C)\), niin induktio-oletuksen nojalla
\(v(B) = v(C) = 1\). Implikaation totuusmääritelmän mukaan tällöin
\(v(A) = 1\).

Siis kaikki propositiosymbolista \(p_0\) sekä konnektiivista \(\to\)
muodostetut propositiolauseet \(A\) toteuttavat \(v(A) = 1\).
Koska \(f(1)=0\), mikään tällainen lause ei esitä funktiota \(f\).
Siis \(\set{\to}\) ei ole täydellinen konnektiivijoukko.
\end{solution}

\begin{exercise}
Voiko päättelyjärjestelmä olla
\begin{enumerate}[label=\alph*), nosep]
\item ehyt muttei täydellinen?
\item täydellinen muttei kumoustäydellinen?
\item kumoustäydellinen muttei täydellinen?
\item kumoustäydellinen muttei ehyt?
\end{enumerate}
Perustele. Keksitkö sopivia esimerkkejä?
\end{exercise}

\begin{solution}
\leavevmode\par
\begin{enumerate}[label=\alph*), nosep]
\item Olkoon \(\mathcal{S}\) päättelyjärjestelmä, joka sisältää vain konjunktion
eliminointisäännöt
\[
\frac{A \land B}{A}
\qquad
\frac{A \land B}{B}.
\]
Järjestelmä \(\mathcal{S}\) on ehyt, sillä jos \(v(A \land B) = 1\), niin
\(v(A) = v(B) = 1\). Siten kumpikin päättelysääntö säilyttää totuuden.
Koska jokainen päättelyn askel säilyttää totuuden, kaikki johdetut
lauseet ovat oletusten loogisia seurauksia. Järjestelmä \(\mathcal{S}\)
ei ole täydellinen. Nimittäin \((p_0 \lor p_1)\) on oletuksen \(p_0\)
looginen seuraus, mutta sitä ei voida johtaa järjestelmässä \(\mathcal{S}\).
Oletukseen \(p_0\) ei voida soveltaa kumpaakaan eliminointisääntöä,
joten siitä ei voida johtaa uusia lauseita.
\item Ei voi. Oletetaan, että päättelyjärjestelmä on täydellinen.
Jokaisen ristiriitaisen oletusjoukon looginen seuraus on ristiriita,
sillä mikään totuusjakauma ei toteuta oletuksia. Täydellisyyden nojalla
oletuksista voidaan siis johtaa ristiriita. Näin ollen järjestelmä on
kumoustäydellinen.
\item Voi. Olkoon \(\mathcal{S}\) päättelyjärjestelmä, jonka ainoa sääntö
sallii ristiriidan johtamisen ilman oletuksia. Järjestelmä on kumoustäydellinen,
sillä ristiriita voidaan johtaa jokaisesta ristiriitaisesta oletusjoukosta.
\(\mathcal{S}\) ei kuitenkaan ole täydellinen, sillä esimerkiksi oletuksesta
\(p_0\) ei voida johtaa johtopäätöstä \(p_0 \lor p_1\), vaikka se on
oletuksen looginen seuraus. Ainoa päättelysääntö tuottaa vain ristiriidan,
joten sillä ei voida johtaa lausetta \(p_0 \lor p_1\).
\item Kohdan (c) järjestelmä on kumoustäydellinen, muttei ehyt, sillä
järjestelmällä voidaan johtaa ristiriita myös toteutuvasta oletusjoukosta.
Esimerkiksi oletuksesta \(p_0 \land p_1\) voidaan johtaa ristiriita, vaikka
ristiriita ei ole tämän oletuksen looginen seuraus: oletus on tosi, kun
molemmat propositiosymbolit \(p_0\) ja \(p_1\) saavat totuusarvon \(1\).
\end{enumerate}
\end{solution}

\end{document}
